\section{Computational Cost and Environment}
\label{sec:computeexperiment}

\subsection{Computational Cost Estimates}
\label{sec:computecost}

The MACE framework adds two computationally significant components for processing multimodal datasets: a BERT-based text encoder and a ResNet-based visual encoder within the multimodal content branch, alongside the social-context processing branch inherited from DANES \citep{truica2024danes} and GETAE \citep{GETAE2025}. The proposed methodology also evaluates datasets larger than those used in the original DANES, GETAE, and HMCAN experiments, making evidence-based cost estimation important for assessing feasibility and managing budget and schedule constraints.

For the multimodal content branch, HMCAN \citep{Qian2021} employs BERT\textsubscript{BASE} (110M parameters) as a fine-tuned text encoder and a frozen pre-trained ResNet-50 as the visual encoder. The ResNet-50 contributes a fixed per-sample feature-extraction cost of approximately $3.8 \times 10^{9}$ FLOPs per forward pass \citep{He2016}, while BERT fine-tuning adds a full gradient-update cost; \citet{Devlin2019} report that fine-tuning can be completed in at most a few hours on a single GPU. Caching ResNet-50 and BERT embeddings where experimental conditions permit will substantially reduce repeated computation. These figures suggest that the multimodal content branch is feasible on a single GPU, provided embeddings are cached appropriately.

The social-context branches of DANES and GETAE introduce additional cost through graph construction, propagation modelling, and node-embedding generation \citep{truica2024danes, GETAE2025}. For Stage~1 replication on Twitter15 and Twitter16, graph sizes are modest; this is consistent with GETAE, where experiments were conducted in Google Colab without hardware acceleration \citep{GETAE2025}. Stage~2 dataset adaptation will assess whether larger candidate datasets can be processed without changing the baseline social-context architecture. MuMin \citep{mumindataset}, with approximately 21.5 million tweets and 1.99 million users, may present a substantially greater scalability challenge. \citet{Chiang2019} show that on a graph of comparable scale, full-graph GNN training requires 11.2GB of GPU memory for a 3-layer model and fails from out-of-memory errors at 4 layers, identifying graph scale rather than model size as the binding constraint. If adapting MuMin would require subgraph-based mini-batch training or other substantial changes to the published DANES/GETAE-style social-context branch, those changes will be treated as dataset feasibility findings rather than silently introduced into the main controlled experiment.

Accordingly, computational resource demand will be evaluated as part of dataset selection. FakeNewsNet and MuMin will therefore be assessed during Stage~2 rather than assumed suitable in advance. FakeNewsNet may be more feasible for primary Stage~3 experiments if it provides sufficient image availability and social-context information while allowing the baseline social-context branch to be preserved without major architectural changes. If MuMin or another candidate dataset satisfies the feasibility criteria within the available hardware and schedule, it may be included as an additional cross-dataset generalisation test for RQ3.

\subsection{Recommended Experimental Environment}
\label{sec:environment}

The recommended environment is a single-GPU workstation or cloud instance with at least 24GB GPU VRAM, 128GB system RAM, 16--32 CPU cores, and 1--2TB of fast local storage for datasets, checkpoints, and cached embeddings. An A100-class GPU is preferred, although lower-memory GPUs may be sufficient for pilot runs if batch size, mixed-precision training, and embedding caching are used. This environment is below the four-GPU NVIDIA DGX Station used by DANES \citep{truica2024danes}, but substantially above the unaccelerated Google Colab environment used by GETAE \citep{GETAE2025}. This reflects the additional cost of the multimodal content branch, image feature extraction, and larger Stage~2 dataset screening. Where hardware limits are encountered, the implementation will use mixed-precision training, cached embeddings, reduced batch sizes, and dataset subsampling for pilot runs. Compute reporting will follow the MLPerf convention of measuring wall-clock time to a defined quality target \citep{Mattson2020}, and will include hardware configuration, peak GPU memory, peak system RAM, inference time, and storage requirements alongside classification performance.

Based on the recommended experimental environment, the total compute requirement across all three stages is estimated to be in the order of 50--100 GPU hours. This estimate should be treated as indicative rather than fixed, because the final requirement will depend on dataset size, image availability, batch size, early stopping behaviour, and the number of ablation conditions required. The majority of this cost is expected to occur in Stage~3, where the full MACE framework and ablation conditions are evaluated across both DANES and GETAE backbones.

For cloud execution, cost should be reported separately as GPU-hours and instance-hours. A100-class cloud pricing varies substantially by provider, region, reservation model, and whether the instance is single-GPU or multi-GPU. For example, an AWS p4d.24xlarge instance provides eight A100 GPUs rather than a single GPU, so its hourly instance price should not be treated as a single-GPU-hour cost. As a conservative planning figure, the project should allow approximately AUD~\$500--\$1000 for cloud compute, storage, pilot runs, failed runs, and data-transfer overheads. Where university HPC resources are available, direct per-hour costs may be substantially reduced or eliminated; however, the GPU-hour estimate remains useful for scheduling, queue-time planning, and reproducibility reporting.

\subsection{Evaluation Matrix Template}
\label{sec:evalmatrix}

\begin{table}[htbp]
\centering
\small
\renewcommand{\arraystretch}{1.2}

\begin{tabularx}{\textwidth}{|X|c|c|c|c|c|c|}
\hline
\textbf{Experimental Condition} &
\textbf{\begin{tabular}[c]{@{}c@{}}Train Time\\ (hrs)\end{tabular}} &
\textbf{\begin{tabular}[c]{@{}c@{}}Inference\\ (ms/sample)\end{tabular}} &
\textbf{\begin{tabular}[c]{@{}c@{}}Peak GPU\\ (GB)\end{tabular}} &
\textbf{\begin{tabular}[c]{@{}c@{}}Peak RAM\\ (GB)\end{tabular}} &
\textbf{\begin{tabular}[c]{@{}c@{}}Storage\\ (GB)\end{tabular}} &
\textbf{FLOPs} \\ \hline

\textbf{Stage 3A:} Dataset Baseline & & & & & & \\ \hline
\textbf{Stage 3B:} Ablation 1 (Text-only) & & & & & & \\ \hline
\textbf{Stage 3C:} Ablation 2 (Multimodal) & & & & & & \\ \hline
\textbf{Stage 3D:} Full MACE Framework & & & & & & \\ \hline

\end{tabularx}

\caption{Evaluation matrix template for computational efficiency cost across Stage 3 conditions.}
\label{tab:efficiency_matrix}
\end{table}